\documentclass[final]{beamer}

\mode<presentation> { \usetheme{Rochester} }
\setbeamersize{text margin left=0.6cm,text margin right=0.6cm}
\setbeamertemplate{navigation symbols}{}
\setbeamertemplate{headline}{
  \leavevmode
  \begin{beamercolorbox}[wd=\paperwidth,dp=2ex]{subsection in head/foot}
    \usebeamerfont{subsection in head/foot}
    \begin{columns}[T]
      \begin{column}{.7\paperwidth}
        \vskip4ex
        \begin{flushright}
          \usebeamercolor{title in headline}{
            \color{fg}\textbf{\huge{\inserttitle}}\\[1ex]
          }
          \usebeamercolor{author in headline}{
            \color{fg}\Large{\insertauthor}\\[1ex]
          }
          \usebeamercolor{institute in headline}{
            \color{fg}\Large{\insertinstitute}\\[1ex]
          }
        \end{flushright}
      \end{column}
      \begin{column}{.25\paperwidth}
        \vskip2.0cm
        \begin{flushright}
          \includegraphics[width=.95\linewidth]{figures/TU_Logo_white.png}
        \end{flushright}
      \end{column}
    \end{columns}
    \vskip0cm
  \end{beamercolorbox}
%  \begin{beamercolorbox}[wd=\paperwidth]{section in head/foot}
%    \rule{0pt}{2pt}
%  \end{beamercolorbox}
}

\setbeamertemplate{footline}{
  \leavevmode
  \begin{beamercolorbox}[leftskip=1cm,rightskip=1cm]{subsection in head/foot}
    \usebeamerfont{section in head/foot}
    \vskip1ex
    \normalsize
    \insertinstitute
    \hfill
    \texttt{\correspondingAuthorEmail}
    \hfill
    \texttt{\website}
  \end{beamercolorbox}
  %\begin{beamercolorbox}[wd=\paperwidth]{section in head/foot}
    %\rule{0pt}{2pt}
  %\end{beamercolorbox}
}



\usepackage[english]{babel}
\usepackage[utf8]{inputenc}
\usepackage{amsmath,amsthm,amssymb,latexsym}
\usepackage{multirow}
\usepackage[orientation=portrait,size=a0,scale=1.4]{beamerposter}
\usepackage{tikz}
\usetikzlibrary{arrows}

\usepackage[
  hyperref=false,
  url=false,
  abbreviate=true,
  maxnames=1,
  date=year,
  style=nature,
  backend=bibtex,
  doi=false
]{biblatex}
\AtEveryBibitem{\clearfield{note}}
\AtEveryBibitem{\clearfield{title}}
\bibliography{main.bib}


\title{
  Equation of motion coupled cluster theory using periodic boundary conditions
}
\author{\underline{Alejandro Gallo}, Felix Hummel and Andreas Gr\"uneis}
\def\correspondingAuthorEmail{alejandro.gallo@tuwien.ac.at}
\def\website{http://cqc.tuwien.ac.at}
\institute[]{
  Technical University of Vienna
}


\begin{document}
\nocite{*}
\begin{frame}[fragile]{}

  \begin{block}{\large Abstract}
    We present recent progress in the implementation of an equation of motion
    coupled cluster theory method using periodic boundary conditions and a plane
    wave basis set.  We seek to study electronic ground and excited states in
    Nitrogen Vacancy centers in diamond which have become a promising candidate
    for a bulk room temperature quantum information processing device. Technical
    details of the implementation and preliminary results will be presented.
  \end{block}
  \begin{columns}[t]
    \begin{column}{.49\linewidth}
      \begin{block}{\large Computational details}
        \begin{center}
\begin{tikzpicture}[style={
    blob/.style={
      node distance=4cm,
      minimum size=2cm,
    },
    cc4sstyle/.style={
      blob, circle,
      line width=5,
      inner color=green,
      outer color=red,
      draw,
    },
    ctfstyle/.style={
      blob, circle,
      line width=5,
      fill=green,
      draw=green!40!black,
    },
    vaspstyle/.style={
      blob, circle,
      line width=5,
      fill=red,
      draw=red!40!black,
    },
    myarrow/.style={->, shorten >=10pt, line width=3pt}
  }]
  \draw[]
    node[cc4sstyle] (cc4s) {\texttt{cc4s}}
    node[vaspstyle, below of=cc4s, left of=cc4s, yshift=-2cm] (vasp) {
      \texttt{VASP}
    }
    node[ctfstyle, below of=cc4s, right of=cc4s, yshift=-2cm] (ctf) {
      \texttt{CTF}
    }
  ;
  \def\vaspscale{0.7}
  \draw[myarrow]
    (vasp.north) |- (cc4s.west)
    node[pos=0.4, left, scale=\vaspscale] (planewaves) {
      $\displaystyle
      \psi_{p}(\mathbf{r}) =
      \sum_{\mathbf{G}}
        \psi_{G}^{p}
        e^{i \mathbf{r} \cdot \mathbf{G}}
      $
    }
    node[below of=planewaves, xshift=-3cm, node distance=2cm, scale=\vaspscale]
    (gamma) {
      $\displaystyle
      \Gamma^{p}_{q\mathbf{G}} =
      \sqrt{
        \frac{4\pi}{|\mathbf{G}|^{2}}
      }
      \int\mathrm{d}\mathbf{r}
        \psi_{p}^{*}(\mathbf{r})
        \psi_{q}(\mathbf{r})
        e^{-i\mathbf{G} \cdot \mathbf{r}}
      $
    }
    node[below of=gamma, node distance=3cm, scale=\vaspscale]
    (vpqrs) {
      $\displaystyle
      V^{pq}_{rs}
      =
      \sum_{\mathbf{G}}
        \Gamma^{p\mathbf{G}}_{q}
        \Gamma^{r}_{s\mathbf{G}}
      $
    }
  ;
  \def\ctfscale{\vaspscale}
  \draw[myarrow]
    (ctf.north) |- (cc4s.east)
    node[pos=0.4, scale=\ctfscale, right] (ctftv) {
      \texttt{
        E[""] =
        T
        ["{\color{red}a}{\color{blue}b}{\color{green!50!black}i}{\color{brown}j}"] *
        V
        ["{\color{green!50!black}i}{\color{brown}j}{\color{red}a}{\color{blue}b}"]
      }
    }
    node[below of=ctftv, xshift=0.2cm, node distance=2cm, scale=\ctfscale] {
      \texttt{
        V
        ["{\color{red}i}{\color{blue}j}{\color{red}a}{\color{blue}b}"]
        =
        $\Gamma$["{\color{red}ia}{\color{yellow!20!green}G}"]
        *
        $\Gamma$["{\color{blue}jb}{\color{yellow!20!green}G}"]
      }
    }
  ;
\end{tikzpicture}
\end{center}

      \end{block}
      \begin{block}{\large Coupled-Cluster Ansatz}
        \begin{itemize}
\item
  %
  \[
    \hat{T} =
    \underbrace{
    \sum_{a,i}^{}
      {\color{blue}t^{a}_{i} }
      \hat{a}^{\dagger}_{a} \hat{a}_{i}
    +
    \sum_{a,b,i,j}
      {\color{blue}t^{ab}_{ij} }
      \hat{a}^{\dagger}_{a} \hat{a}^{\dagger}_{b}
      \hat{a}_{j}           \hat{a}_{i}
    }_{\mathrm{CCSD}}
    + \cdots
    =
    \tikz[baseline]\node{\includegraphics[scale=2]{figures/diagrams/t1.pdf}};
    +
    \tikz[baseline]\node{\includegraphics[scale=2]{figures/diagrams/t2.pdf}};
    +
    \cdots
  \]
  %
  Exponential \textit{Ansatz} for dynamic correlation
  %
  \[
    \left | \mathrm{CC} \right \rangle  =
    e^{\hat{T}} \left | 0 \right \rangle
    =
    1 +
    \tikz[baseline]\node{\includegraphics[scale=2]{figures/diagrams/t1.pdf}};
    +
    \tikz[baseline]\node{\includegraphics[scale=2]{figures/diagrams/t2.pdf}};
    +
    \tikz[baseline]\node{\includegraphics[scale=2]{figures/diagrams/t1.pdf}};
    \tikz[baseline]\node{\includegraphics[scale=2]{figures/diagrams/t1.pdf}};
    +
    \tikz[baseline]\node{\includegraphics[scale=2]{figures/diagrams/t1.pdf}};
    \tikz[baseline]\node{\includegraphics[scale=2]{figures/diagrams/t2.pdf}};
    +
    \cdots
  \]
\item
  Size consistency and extensivity
\end{itemize}


      \end{block}
      \begin{block}{\large Coupled-Cluster Ansatz}
        \begin{itemize}
  \item
    In one sentence:
    \begin{center}
      \textit{
      Perform \emph{CI} on the similarity-transformed Hamiltonian $\bar{\mathsf{H}}$
      instead of $H$.
    }
    \end{center}
  \item
    One striking consequence: for EOM-CCSD up to 6 body hamiltonian,
  $$
    \bra{^{abcdef}_{ijklmn}} \bar{\mathsf{H}} \ket{0}
    =
    \begin{tikzpicture}[baseline, line width=2, scale=5]
      \def\angle{15}
\def\tLength{0.5}
\def\tHeight{0.5}
\newcommand{\tamplitude}[1]{
  % define points
  \draw[draw=none]
    (0,0)
      coordinate (#1bottoml)
    ++ (\tLength,0)
      coordinate (#1bottomr)
    (#1bottoml) ++ (90 - \angle:\tHeight)
      coordinate (#1topparticlel)
    (#1bottoml) ++ (90 + \angle:\tHeight)
      coordinate (#1topholel)
    (#1bottomr) ++ (90 - \angle:\tHeight)
      coordinate (#1topparticler)
    (#1bottomr) ++ (90 + \angle:\tHeight)
      coordinate (#1topholer)
  ;

  \draw[] (#1bottoml) -- (#1bottomr);

}

\draw[dashed]
  (0,0) coordinate (vl) --++ (0.7, 0) coordinate (vr)
  node [pos=0.5, rectangle, scale=0.9, above] {
    $V^{\ ij}_{ab}$
  }
;

% t amplitude left
\begin{scope}[yshift=-25, xshift=-15]
  \tamplitude{leftamplitude}
\end{scope}
\draw[-<-] (leftamplitudebottomr) -- (vl);
\draw[->-] (leftamplitudebottoml) -- (leftamplitudetopparticlel);
\draw[-<-] (leftamplitudebottoml) -- (leftamplitudetopholel);
\draw[->-] (leftamplitudebottomr) -- (leftamplitudetopholer);

% t amplitude leftleft
\begin{scope}[yshift=-15, xshift=-45]
  \tamplitude{leftleftamplitude}
\end{scope}
\draw[->-] (leftleftamplitudebottomr) -- (vl);
\draw[->-] (leftleftamplitudebottoml) -- (leftleftamplitudetopparticlel);
\draw[-<-] (leftleftamplitudebottoml) -- (leftleftamplitudetopholel);
\draw[-<-] (leftleftamplitudebottomr) -- (leftleftamplitudetopholer);

% t amplitude right
\begin{scope}[yshift=-25, xshift=35, xscale=-1]
  \tamplitude{rightamplitude}
\end{scope}
\draw[-<-] (rightamplitudebottomr) -- (vr);
\draw[->-] (rightamplitudebottoml) -- (rightamplitudetopparticlel);
\draw[-<-] (rightamplitudebottoml) -- (rightamplitudetopholel);
\draw[->-] (rightamplitudebottomr) -- (rightamplitudetopholer);

% t amplitude rightright
\begin{scope}[yshift=-15, xshift=65, xscale=-1]
  \tamplitude{rightrightamplitude}
\end{scope}
\draw[->-] (rightrightamplitudebottomr) -- (vr);
\draw[->-] (rightrightamplitudebottoml) -- (rightrightamplitudetopparticlel);
\draw[-<-] (rightrightamplitudebottoml) -- (rightrightamplitudetopholel);
\draw[-<-] (rightrightamplitudebottomr) -- (rightrightamplitudetopholer);

    \end{tikzpicture}
  $$

  \item
    Linear (CI) \textit{Ansatz} on top of ground state coupled-cluster.
    \begin{align*}
      \ket{r_i} = \hat{R} ^{(i)} e^{\hat{T}} \ket{0}, \qquad
      \bra{l_i} = \bra{0} e^{-\hat{T}} \hat{L}^{(i)}
    \end{align*}
    where
      \begin{align*}
        \hat{R} ^{(i)} &= r _{0}
         + r ^{a} _{i} \hat{c} ^{\dagger} _{a} \hat{c} _{i}
         + r ^{ab} _{ij} \hat{c} ^{\dagger} _{a} \hat{c} ^{\dagger} _{b}
                \hat{c} _{j} \hat{c} _{i}
         + \ldots
        =
        \tikz[baseline]\node{\includegraphics[scale=2]{figures/diagrams/r1.pdf}};
        +
        \tikz[baseline]\node{\includegraphics[scale=2]{figures/diagrams/r2.pdf}};
        +
        \cdots
        \\
        \hat{L} ^{(i)} &= l _{0}
                       + l ^{i} _{a} \hat{c} _{i}^{\dagger} \hat{c} _{a}
                       + l ^{ij} _{ab}
                          \hat{c} _{j} ^{\dagger} \hat{c} _{i}^{\dagger}
                          \hat{c} _{a} \hat{c} _{b}
                       + \ldots
        =
        \tikz[baseline]\node{\includegraphics[scale=2,angle=180]{figures/diagrams/r1.pdf}};
        +
        \tikz[baseline]\node{\includegraphics[scale=2,angle=180]{figures/diagrams/r2.pdf}};
        +
        \cdots
      \end{align*}

  \item
    Non-hermitian diagonalization problem of
      $ \bar{\mathsf{H}} = e^{-\hat{T}} \hat{H} e^{\hat{T}}$

  \item
    Left and right eigenvectors $ \Rightarrow $ accessibility to calculate
    properties

  \item
    Contained size-extensivity error

  \item
    Davidson-based diagonalization procedure

\end{itemize}

%\textbf{Davidson-based diagonalization procedure}
%\begin{itemize}
%
  %\item
    %Get initial \textit{right} basis
    %$ B = \left \{ b_{1}, \ldots, b_{k} \right \} $
    %to obtain $ k $ excited states
%
  %\item
    %Compute reduced matrix and diagonalize
    %$
      %h_{ij} =
      %\bra{b_{i}} \bar{\mathsf{H}} \ket{b_{j}}
    %$
%
  %\item
    %Get $ k $ lowest eigenvalues and eigenvectors and compute corrections
    %for them
%
  %\item
    %If the corrections are significant, add them to the basis
%
  %\item
    %Basis refreshment: at some iterations, the basis gets collapsed
    %to have only $ k $ vectors
%
%\end{itemize}

      \end{block}
    \end{column}

    % -------------------------------------------------------------

    \begin{column}{.49\linewidth}

      \begin{block}{\large Color centers in alkali halides}
        \includegraphics[width=0.8\textwidth]{figures/fcenters/plot.pdf}
      \end{block}

      \begin{block}{\large Molecule benchmark example: Formaldehyde (CH$_2$O)}
          %\begin{columns}[t]
            %\begin{column}{.2\linewidth}
              %\begin{itemize}
                %\item
                  %Calculation with pseudized-gaussians \cite{10.10631.4961301}
              %\end{itemize}
            %\end{column}
            %\begin{column}{.8\linewidth}
            %\end{column}
          %\end{columns}
          \includegraphics[width=\textwidth]{figures/formaldehyde/plot.pdf}
          \begin{itemize}
            \item
              Calculation with pseudized-gaussians
                (see \cite{10.10631.4961301}).
            \item
              Singlet and triplet states retrieved.
          \end{itemize}
      \end{block}

      \begin{block}{\large References \& Acknowledgements}
        \printbibliography
      \end{block}

    \end{column}

  \end{columns}

\end{frame}
\end{document}
